\section{Periodic trace presentations of the common middle family}

\begin{theorem}[A common polynomial family with periodic trace presentations]
    \label{thm:physical_aeval_periodic_middle_family}
    \lean{TNLean.MPS.AreaLaw.exists_common_physical_aeval_periodic_family_fin}
    \uses{thm:physical_aeval_fin_middle_family,thm:trace_chain_closing_absorption}
    Retain the hypotheses (\ref{eq:physical_polynomial_family_hypotheses}),
    the supplied window operators, core tensors and closing matrix of
    (\ref{eq:physical_polynomial_supplied_core}), and the notation
    (\ref{eq:physical_polynomial_two_bounds}). Write $\beta=\max(1,B)$.
    One may choose $N\le D^2$ and matrices $W_0,\ldots,W_{N-1}$ on
    $\mathcal C(J)$ with all the properties of
    Theorem~\ref{thm:physical_aeval_fin_middle_family} and the following
    additional property. For each $r$, there are a nonnegative integer
    $\chi_r\le\beta^k$ and site-dependent matrices
    $T_{r,v}(x,y)\in\C^{\chi_r\times\chi_r}$ such that
    \begin{align}
        W_r(\sigma,\tau)
        &=\operatorname{tr}\left(
        \prod_{v=0}^{|J|-1}T_{r,v}(\sigma_{j_v},\tau_{j_v})\right).
        \label{eq:physical_periodic_middle_trace}
    \end{align}
    Here $j_0<\cdots<j_{|J|-1}$ enumerate $J$, and
    $\sigma,\tau\in\mathcal C(J)$.
    In particular, the same matrices $W_r$ retain their open-boundary bond
    bound $\mathcal B$ and their presentations
    (\ref{eq:physical_polynomial_trace_presentation}) with a closing matrix.
    The number, matrices and both trace presentations are chosen before
    every supported exterior pair $P_L,P_R$ and every polynomial
    $P\in\C[X]$ of natural degree at most $k$.
    For each such pair and polynomial there are exterior matrices $M_r$ with
    \begin{align}
        P(H)&=\left(\sum_{r=0}^{N-1}M_r\otimes W_r\right)^{\mathrm{phys}},
        \label{eq:physical_periodic_middle_decomposition}
    \end{align}
    where $H$ is given by (\ref{eq:physical_polynomial_hamiltonian}).
    The zero polynomial and zero supplied virtual dimension are allowed.
\end{theorem}
\begin{proof}
    \uses{thm:physical_aeval_fin_middle_family,thm:trace_chain_closing_absorption}
    Choose the common family and, for each member, its closing matrix and
    tensors from (\ref{eq:physical_polynomial_trace_presentation}).
    The geometric assumptions imply $|J|=(b+\ell+1)-(a+\ell)\ge1$.
    Absorb each closing matrix into the first tensor by
    Theorem~\ref{thm:trace_chain_closing_absorption}. The virtual dimension
    $\chi_r$ and physical matrix $W_r$ are unchanged. Thus the count, the
    open-boundary estimate and (\ref{eq:physical_periodic_middle_decomposition})
    hold for the same family.
\end{proof}

These explicit trace presentations give supporting algebra for
\cite[Proposition~5.7]{Arad2016RigorousRG}.
